\documentclass[12pt]{article}
\usepackage[utf8]{inputenc}
\usepackage{lipsum} % to generate some filler text
\usepackage{fullpage}
\usepackage{bm}
\usepackage{amsmath,amssymb,dsfont, mathtools}
\usepackage[pdftex]{graphicx}
\usepackage{epstopdf}
\usepackage{subcaption}
\usepackage{natbib}
\usepackage{hyperref}
\usepackage{placeins}
\usepackage{algorithm}
\usepackage{algorithmic}
\usepackage{amsmath}
\usepackage{booktabs}


% import Eq and Section references from the main manuscript where needed
% \usepackage{xr}
% \externaldocument{manuscript}
\DeclareMathOperator{\plim}{plim}
\DeclareMathOperator{\vect}{vec}
\DeclareMathOperator{\D}{D}
\DeclareMathOperator{\Diag}{Diag}
\DeclareMathOperator{\argmin}{argmin} 
\DeclareMathOperator{\argmax}{argmax} 
\newcommand{\bs}{\boldsymbol}
\newcommand{\E}{\mathbb{E}}
\newcommand{\tr}{\mathrm{tr}}
\newcommand{\abs}{\mathrm{abs}}
\newcommand{\var}{\mathrm{Var}}
\newcommand{\mb}{\mathbb}
\newcommand{\mf}{\mathbf}
\newcommand{\mc}{\mathcal}
\newcommand{\mk}{\mathfrak}
\renewcommand{\vec}{\text{vec}}
\newcommand{\red}[1]{{\color{red}#1}} 
\newcommand{\h}{\mathbf{h}}
\newcommand{\Hh}{\mathbf{H}}
\newcommand{\W}{\mathbf{W}}
\newcommand{\M}{\mathbf{M}}
\newcommand{\R}{\mathbf{R}}
\newcommand{\U}{\mathbf{U}}
\newcommand{\V}{\mathbf{V}}
\newcommand{\I}{\mathbf{I}}
\newcommand{\Y}{\mathbf{Y}}
\newcommand{\X}{\mathbf{X}}
\newcommand{\Ss}{\mathbf{S}}
\newcommand{\A}{\mathbf{A}}
\newcommand{\J}{\mathbf{J}}
\newcommand{\K}{\mathbf{K}}
\newcommand{\Pp}{\mathbf{P}}
\newcommand{\0}{\mathbf{0}}
\newcommand{\C}{\mathbf{c}}
\newcommand{\f}{\mathbf{f}}
\newcommand{\Lag}{\mathbf{L}}
\newcommand{\B}{\mathbf{B}}
\newcommand{\s}{\mathbf{s}}
\newcommand{\Z}{\mathbf{Z}}
%Bold greek letter using the first two letters
\newcommand{\be}{\boldsymbol{\beta}}
\newcommand{\la}{\boldsymbol{\lambda}}
\newcommand{\rh}{\boldsymbol{\rho}}
\newcommand{\et}{\boldsymbol{\eta}}
\newcommand{\mm}{\boldsymbol{\mu}}
\newcommand{\de}{\boldsymbol{\delta}}
\newcommand{\xx}{\boldsymbol{\xi}}
\newcommand{\bb}{\boldsymbol{\beta}}
\newcommand{\Om}{\boldsymbol{\Omega}}
\newcommand{\Si}{\boldsymbol{\Sigma}}
\newcommand{\Up}{\boldsymbol{\Upsilon}}
\newcommand{\Th}{\boldsymbol{\theta}}
\newcommand{\Ps}{\boldsymbol{\Psi}}
\newcommand{\Ph}{\boldsymbol{\Phi}}
\newcommand{\La}{\boldsymbol{\Lambda}}
\newtheorem{rmk}{Remark}

\newcommand{\xvec}{\boldsymbol}
\newcommand{\xmat}{\mathbf}
\newcommand{\xset}{\mathds}

\newcommand{\y}{\xvec{Y}}
%\newcommand{\K}{\mathbf{K}}
\newcommand{\x}{\xvec{x}}
\newcommand{\z}{\xvec{z}}
\DeclareMathOperator{\e}{\varepsilon}
\usepackage[shortlabels]{enumitem}
\renewcommand{\theenumi}{\Alph{enumi}}

\newcommand{\blue}{\textcolor{blue}}

% package needed for optional arguments
\usepackage{xifthen}
\usepackage{float}
% define counters for reviewers and their points
\newcounter{reviewer}
\setcounter{reviewer}{0}
\newcounter{point}[reviewer]
\setcounter{point}{0}

% This refines the format of how the reviewer/point reference will appear.
\renewcommand{\thepoint}{P\,\thereviewer.\arabic{point}} 

% command declarations for reviewer points and our responses
\newcommand{\reviewersection}{\stepcounter{reviewer} \bigskip \hrule
                  \section*{}}

\newenvironment{point}
   {\refstepcounter{point} \bigskip \noindent {\textbf{Point~\thepoint} } ---\ }
   {\par }

\newcommand{\shortpoint}[1]{\refstepcounter{point}  \bigskip \noindent 
	{\textbf{Reviewer~Point~\thepoint} } ---~#1\par }

\newenvironment{reply}
   {\medskip \noindent \begin{sf}\color{Blue}\textbf{Reply}:\  }
   {\medskip \end{sf}}

\newenvironment{replycont}
   {\medskip \noindent \begin{sf}\color{Blue}\  }
   {\medskip \end{sf}}


 \newenvironment{changes}
   {\medskip \noindent \begin{sf}\color{ForestGreen}\textbf{Changes}:\  }
   {\medskip \end{sf}}

\newcommand{\shortreply}[2][]{\medskip \noindent \begin{sf}\textbf{Reply}:\  #2
	\ifthenelse{\equal{#1}{}}{}{ \hfill \footnotesize (#1)}%
	\medskip \end{sf}}

\usepackage[pagewise, mathlines]{lineno} % Needed to include, otherwise, doesn't work
\usepackage[dvipsnames]{xcolor}

\begin{document}

\author{ }
\bibliographystyle{agsm}
\date{}
\title{Responses to reviewers}
\maketitle

\vspace*{-1.0cm}

\begin{sf}
\noindent \noindent We thank the reviewer for the insightful suggestions provided during this revision round. We have addressed each of the points raised; a detailed point-by-point response is provided below. 

\noindent The revised manuscript incorporates several major enhancements to address the concerns raised and to improve the overall transparency of the approach:

\begin{enumerate}
 \item A \textbf{revised algorithm} where offset levels are implemented by means of a Generalised Least Squares procedure rather than a mean function.
    \item \textbf{A new simulation experiment} specifically designed to demonstrate the precision  of the proposed approximation for each subcomponent.
    \item \textbf{Updated experimental designs} for both the synthetic and real-world data experiments, revised in accordance with the reviewer’s recommendations to ensure greater robustness.
    \item \textbf{A dedicated sensitivity analysis} investigating the impact of different ordering strategies on approximation accuracy.
    \item \textbf{Enhanced technical transparency} regarding the positive semi-definiteness of the implied covariance structure, the details of the optimisaton routines, the selection of comparative benchmarks, and the explicit formulation of key model components such as $H$, $Z$, and $\rho(s)$.
\item Notice that structure also has been reorganized. Now the paper follows: \begin{itemize}
    \item \textbf{Section 1} = Motivation and gap
    \item \textbf{Section 2} = Model + methodological innovation + scalable algorithm
    \item \textbf{Sections 3.1--3.2} = Controlled validation
    \item \textbf{Section 3.3} = Real-data demonstration
    \item \textbf{Section 4}= Conclusion
    \item \textbf{Appendix A--F} = Technical support and supplementary evidence
\end{itemize}

    
\end{enumerate}


\end{sf}






\begin{point}
The core methodological innovation relies on decomposing the covariance
structure and applying the Vecchia approximation independently
to the latent sub-components (e.g., ΣL and Σδ), then reconstructing the
required inverse/log-determinant through the Woodbury identity. The
manuscript currently asserts the validity of this strategy without sufficiently
demonstrating that approximation error in the sub-components
does not propagate destructively to the reconstructed K−1 or to predictive
uncertainty. I recommend including a targeted experiment on
a manageable dataset size where exact GP inference is feasible, comparing
(i) the proposed “Decomposed Vecchia” approach against (ii) an
exact GP benchmark and/or (iii) a “Full Vecchia” approximation applied
directly to the joint model (if implementable). This would provide concrete
empirical assurance that the decomposition strategy is statistically
sound.
\end{point}

\begin{reply}
We thank the reviewer for highlighting these important points, which directly concerns the validity of the proposed decomposition strategy. In response, we have added a new section that explicitly investigates how approximation errors propagate when the Vecchia approximation is applied separately to the latent subcomponents.

Specifically, we conduct a targeted validation experiment in which the decomposed Vecchia approximation is compared with exact multi-fidelity Gaussian process (MFGP) inference on a dataset of manageable size where the exact likelihood can be computed. Table 1 reports the results of this experiment, analysing the behaviour of the approximation under different conditioning strategies (Nearest Neighbour and Correlation-based conditioning) and for varying neighbour sizes. The analysis evaluates several key quantities involved in likelihood reconstruction, including $K^{-1}y$, $y^\top K^{-1}y$, and $\log |K|$, together with predictive accuracy metrics. These diagnostics allow us to assess directly whether approximation errors in the latent components propagate destructively when reconstructing the joint covariance structure.

In addition, in the synthetic simulation experiment we now report the performance of the exact MFGP model alongside the approximated version. This comparison provides a direct empirical assessment of the approximation error in a controlled setting where the ground-truth model is known.

To further address the reviewer’s concern regarding predictive uncertainty, we also evaluate uncertainty calibration through prediction-interval coverage in the real-data experiment (Section 3.3). This complementary analysis confirms that the proposed decomposed Vecchia approximation preserves not only the accuracy of the reconstructed likelihood quantities but also the reliability of the resulting predictive uncertainty estimates.

Regarding the suggestion of applying a fully stacked Vecchia approximation to the joint multi-fidelity process, we emphasise that such an approach requires particular care because the conditioning scheme must respect the cross-fidelity dependence structure. A naive construction obtained by stacking $[X_L, X_H]$ and applying a standard Vecchia approximation does not automatically preserve the intended cross-covariance relationships. In practice, several methodological challenges arise.

\paragraph{Fidelity-aware conditioning.}
Standard geometric neighbour selection is agnostic to fidelity labels and to the cross-correlation parameter $\rho$. Consequently, the resulting conditioning sets may not reflect the dependence structure implied by the multi-fidelity kernel. For example, when $\rho$ is small, cross-fidelity neighbours may contribute negligible information while still occupying conditioning slots, thereby reducing approximation quality.

\paragraph{Ordering sensitivity and asymmetry.}
Vecchia approximations are inherently directional because each observation conditions only on previously ordered points. In a stacked setting, the ordering of low- and high-fidelity observations directly determines the direction of conditioning. For instance, if the data are ordered as $[L_1, L_2, H_1, H_2]$, high-fidelity points may condition on low-fidelity ones but not vice versa; reversing the order produces the opposite effect. Although the true joint covariance is symmetric, the induced conditional structure becomes asymmetric, which may lead to different approximations of the cross-correlation parameter $\rho$ under different orderings.

\paragraph{Kernel–conditioning coherence.}
The interaction between the multi-fidelity kernel parameters (including $\rho$ and potentially distinct spatial–temporal lengthscales) and local conditional regressions is nontrivial. Geometric neighbour selection does not necessarily align with the strongest correlations implied by the kernel, particularly when the lengthscales differ across fidelities.

\paragraph{Stability of hyperparameter estimation.}
Because cross-fidelity dependence is approximated locally through neighbour sets, the estimation of $\rho$ and related parameters may become sensitive to the composition and ordering of the conditioning sets, potentially affecting optimisation stability.

Addressing these issues would require a dedicated theoretical and computational investigation to ensure that the resulting approximation faithfully represents the intended multi-fidelity structure.

\paragraph{Proposed approach.}
The approach adopted in this work avoids these difficulties by construction. The low- and high-fidelity processes are approximated separately at the Vecchia level, so neighbour selection occurs strictly within each fidelity. Cross-fidelity dependence is subsequently introduced through a structured coupling mechanism rather than through mixed conditioning sets.

This separation
\begin{itemize}
\item eliminates fidelity-blind neighbour selection,
\item removes ordering-induced asymmetry across fidelities,
\item preserves coherence between the kernel structure and the conditioning scheme, and
\item improves numerical stability in hyperparameter estimation.
\end{itemize}

For these reasons, the proposed formulation provides a principled and computationally stable framework for multi-fidelity modelling. A fully stacked Vecchia construction with fidelity-consistent conditioning remains an interesting direction for future research.
\end{reply}


\begin{point}
The computational efficiency claim (e.g., near-linear scaling) relies heavily
on the matrix $H = Σ−1
w +A⊤D−1 ϵ$ A being “sparse and banded.” However,
the sparsity pattern of H is not intrinsic; it can depend substantially
on the ordering of the joint vector y (LF vs. HF ordering, space–time
ordering, etc.) and on the neighbour construction used in the Vecchia
approximation. In particular, interactions introduced through A can destroy
banded structure under certain orderings, potentially negating the
claimed complexity. The authors should therefore (i) state explicitly the
ordering strategy and neighbour construction used in all experiments, and
(ii) provide either a mathematical justification for sparsity/bandedness
under that ordering or clear empirical evidence (e.g., sparsity visualizations
/ nnz counts) under several plausible orderings (e.g., LF-first/HFlast
versus interleaved or space–time orderings
   
\end{point}

\begin{reply}
We thank the reviewer for highlighting this point. As noted, the ordering strategy significantly impacts both the sparsity and the banded structure of the matrix. We have revised Section 2.2 to reflect this more accurately. In all experiments, observations are ordered using a space-major spatio-temporal ordering combined with an Approximate Minimum Degree (AMD) permutation to further reduce fill-in and preserve sparsity in the resulting matrices.
Furthermore, we have added a new section to Appendix B that analyzes non-zero ($nnH$) counts and the degree of sparsity under various ordering strategies. This appendix also evaluates the accuracy of the approximation and includes figures illustrating the banded structure and the sparsity patterns resulting from the Approximate Minimum Degree (AMD) permutation used in our algorithm.
\end{reply}

\begin{point}
In the section where the constant rescaling parameter $\rho$ is replaced by
a spatially varying function $\rho(s)$, the manuscript does not clearly and
explicitly state the resulting functional form of the HF–HF covariance
block and the HF–LF cross-covariance. Because covariance kernels must
be symmetric positive semidefinite (PSD), it is important to confirm the
model’s validity under the chosen ρ(s) formulation. Please explicitly
write down the mathematical form used for the HF–HF block kMF (s, s′)
and for the HF–LF cross-covariance, and include a brief justification
(e.g., a short lemma/remark using the hierarchical representation and the
implied block covariance factorization) that the resulting covariance is symmetric and PSD. Right now it is unclear whether the implementation
uses a PSD-preserving form such as ρ(s)ρ(s′)kL(s, s′) (or the equivalent
matrix congruence form) or another formulation.
\end{point}

\begin{reply}

We thank the reviewer for highlighting the importance of the kernel's properties. We have added \textbf{Appendix E} to the manuscript, which explicitly defines the kernel functions and provides formal proofs of their \textbf{positive semidefiniteness (PSD)}. 

We confirm that the implementation utilizes the PSD-preserving form:
\begin{equation*}
    k_{HH}(s, s') = \rho(s)\rho(s')k_L(s, s') + k_\delta(s, s')
\end{equation*}
This formulation is theoretically justified by representing the model as a linear transformation of latent Gaussian processes. Consequently, the resulting covariance is a \textbf{congruence transformation} of a PSD matrix, which inherently guarantees the PSD property.
\end{reply}


\begin{point}
    The manuscript acknowledges that a linear ρ(s) model can yield negative
rescaling coefficients, which can in turn generate physically impossible
negative wind speed predictions. The paper mentions “warping” as a
possible solution, but treats it as optional. For an environmental modelling
application, systematic violation of basic physical constraints is
a substantial drawback. I strongly recommend either integrating the warping/transformation into the main pipeline (and showing results) or
providing a quantitative justification for why allowing negative predictions
is acceptable in this context (e.g., negative values occur rarely and
with negligible magnitude, or only occur in regimes where the practical
impact is minimal).
\end{point}

\begin{reply}
We thank the reviewer for this insightful comment regarding the integration of the warping adjustment. After careful consideration, we have decided to maintain the warping as an optional diagnostic or post-processing step rather than integrating it into the core algorithm, for several reasons: Application-Specific Requirement: The occurrence of negative predictions is highly dependent on the physical nature of the data. While negative values are physically implausible for variables like wind speed, they are entirely admissible—and often necessary—for other environmental data, such as temperature (Celsius) or pressure anomalies. Keeping the main algorithm 'un-warped' ensures its generalizability across different environmental datasets. Empirical Frequency: In our extensive testing, the majority of model runs produced physically consistent results without any warping adjustment. We view the occurrence of negative values not as a structural failure of the model, but as a rare edge-case anomaly. Pathology of the Anomaly: Our analysis indicates that these anomalies typically only emerge under a specific convergence of two conditions: (1) observations located extremely close to the zero-boundary, and (2) the presence of a negative rescaling factor ($\rho$). In most operational contexts, this is an infrequent occurrence. In the real data experiment we tested 108 models. The table below illustrate the frequency of occurrence of negative values and their magnitude:

\begin{table}[h]
\centering
\caption{Granular Analysis of Negative Predictions ($y_{pred} < 0$)}
\label{tab:detailed_negatives}
\begin{tabular}{@{}lcc@{}}
\toprule
\textbf{Error Category} & \textbf{No. of Models} & \textbf{Average Minimum Value} \\ \midrule
Models with 0 negatives    & 90 & 0.0000 \\
Models with 1 negative     & 8  & -0.0818 \\
Models with 2 negatives    & 3  & -0.2232 \\
Models with 3 negatives    & 3  & -0.0812 \\
Models with 4 negatives    & 1  & -0.1597 \\
Models with 5 negatives    & 1  & -0.4586 \\
Models with 6 negatives    & 1  & -0.5877 \\
Models with $>6$ negatives & 1  & -6.3333 \\ \bottomrule
\end{tabular}
\vfill
\end{table}


\noindent
It should be noted that, where negative values occur, they are generally of negligible magnitude and do not significantly impact the overall model integrity. The only exception is the \textbf{MFGP} configuration employing \textbf{Constant GLS} and \textbf{Adaptive} $\rho(\mathbf{s})$ trained on \textbf{Station 856}, which produced an unacceptable minimum value of $-6.33$. This outlier underscores the inherent instability of the adaptive $\rho$ framework when implemented without sufficient regularization or positivity constraints. We believe this provides sufficient transparency for users to decide if warping is necessary for their specific application.
\end{reply}




\begin{point}
In the “Clusters” experiment (page 42), the optimised Gaussian process
parameters show extreme variation, with the latitude length scale lϕ
jumping from 0.338 in the full dataset to 7.097 in the clusters. This
degree of variability suggests potential instability/identifiability issues
or overfitting when using a spatially varying ρ(s). Please include a brief
stability diagnostic (e.g., multiple random initializations and/or fold-wise
parameter estimates) and comment on whether $\rho(s)$ is confounded with
kernel hyperparameters or the discrepancy component in this setting.
\end{point}

\begin{reply}
We thank the reviewer for highlighting this important point. In page 42, there is no $ \rho(s) $ as the experiment involves a comparison of simpler Gaussian process only.  However, in the revised version of the paper this section has been removed. As all the experiments were performed from scratch again. The full code for the experiments is provided in GitHub for reproducibility. The Gaussian process models were implemented using the \texttt{fitrgp} function in MATLAB, utilizing an Automatic Relevance Determination (ARD) Squared Exponential kernel. Both the full-scale models (evaluated on a reduced dataset)  and their approximated counterparts failed to yield competitive performance.
\end{reply}


\begin{point}
    Some of the empirical comparisons may be difficult to interpret if the
proposed multi-fidelity model is trained using both HF and LF data
while a baseline (e.g., “GP(yL)”) is trained only on LF observations. In
that case, the observed performance gap may primarily reflect access
to additional HF information rather than an advantage of the proposed
modelling/approximation strategy. Please clarify what training data each
baseline uses and include at least one baseline that has access to HF
information (e.g., an HF-only GP/NNGP, or a simple GP model for HF
that uses LF as a covariate), so that the comparison isolates the benefit
of the multi-fidelity structure.
\end{point}

\begin{reply}
  We appreciate the referee's concern regarding the fairness of the baseline comparisons. We clarify that \textbf{all} GP baselines evaluated in this study—\textbf{GP-L}, \textbf{GP-3D}, and \textbf{GP-4D}—are trained using High-Fidelity (HF) observations as the target variable ($y_H$). The performance differences, therefore, do not arise from unequal access to HF data, but rather from the specific choice of input features used to define the covariance structure. 

To resolve any ambiguity, we have updated the manuscript to precisely define these configurations:
\begin{itemize}
    \item \textbf{GP-L:} $y_H \sim \mathcal{GP}(y_L(\mathbf{x}_H))$, where the model maps the LF signal directly to the HF response ($y_L \rightarrow f \rightarrow y_H$). This baseline tests whether the LF data alone is a sufficient covariate for the HF target.
    \item \textbf{GP-3D:} $y_H \sim \mathcal{GP}(\mathbf{x}_H)$, where the model relies solely on spatio-temporal coordinates ($[\mathbf{s}, t] \rightarrow f \rightarrow y_H$). This represents the standard HF-only GP baseline suggested by the referee.
    \item \textbf{GP-4D:} $y_H \sim \mathcal{GP}(y_L(\mathbf{x}_H), \mathbf{x}_H)$, which jointly exploits the LF information and the spatio-temporal coordinates as features ($[y_L, \mathbf{s}, t] \rightarrow f \rightarrow y_H$).
\end{itemize}

In our preliminary tests on synthetic and reduced real-world datasets, \textbf{GP-3D} proved to be the most competitive baseline. Consequently, we utilized \textbf{GP-3D} for the full-scale real-data experiments to ensure that our multi-fidelity model was compared against the most robust standard GP alternative. We have updated the notation throughout Section 3.2 to ensure this distinction is clear to the reader.
\end{reply}



\begin{point}
    Vecchia/nearest-neighbour approximations can be sensitive to neighbour
size and ordering. The paper would benefit from a concise sensitivity
analysis showing how predictive accuracy, uncertainty calibration (e.g.,
interval coverage), and runtime change with m (and optionally for one
or two alternative orderings). Reporting a small grid would provide
actionable guidance and strengthen the credibility of the computational
claims.
\end{point}

\begin{reply}
We thank the referee for this helpful comment. In response, we added two new experiments to assess the impact of both conditioning and ordering on the proposed approximation. Specifically, Table 1 (In section 3.1) reports a sensitivity analysis of the approximation accuracy under different conditioning choices, while table B.2 presents the corresponding analysis for alternative ordering strategies. 

In addition, we expanded the empirical evaluation in two further ways. First, we now report prediction-interval coverage in the real-data experiment, providing a clearer assessment of uncertainty quantification. Second, in the simulated-data experiment table 2, we included a direct comparison between the approximated and non-approximated algorithms, allowing us to quantify the loss in accuracy induced by the approximation relative to the exact approach.
\end{reply}







\begin{point}
    The paper’s reliance on maximizing the approximate log-likelihood for
parameter estimation is central to the framework, yet the manuscript
provides no details on the numerical optimisaton procedure used. As
this is a computational paper and the likelihood surface for complex
spatial models is often non-convex with local maxima, this omission represents
a significant gap in methodological rigor and reproducibility, particularly
given the parameter instability issues observed in the results.
Please validate the robustness of the parameter estimation by explicitly
reporting the specific optimiser used, the convergence criteria (ϵ), the
initialization strategy for the hyperparameters, and any other related
computation/optimisaton details.
\end{point}

\begin{reply}

The reviewer rightly points out the importance of the numerical optimisaton procedure in ensuring the reliability of the spatial model parameters. To address this, we provide a detailed description of the optimisaton framework used in this study.
In general we used a set of informed initial parameters optimised with a Quasi-Newton (BFGS) procedure.  The set of initial parameters was found with a more complex routine of models trained on a smaller subset of data, more precisely:

\begin{enumerate}
\item*{1. \textbf{optimisaton Procedure}}
Parameter estimation was performed by maximizing the approximate log-likelihood, $\ell(\theta)$, using the \textbf{Quasi-Newton (BFGS)} algorithm. This method was implemented via the \texttt{fminunc} solver in MATLAB. Unlike a standard Newton-Raphson approach which requires the exact Hessian, the BFGS algorithm approximates the Hessian using gradient information, providing a robust balance between convergence speed and computational efficiency for high-dimensional spatial models.

\item{2. \textbf{Parameter Vector and Stability}}
To ensure numerical stability and maintain physical constraints (such as the positivity of variances and length-scales), the hyperparameter vector $\theta$ was optimised in the log-domain. For the most complex model (\textit{GP-scaled $\rho$}), the vector $\theta \in \mathbb{R}^{14}$ is defined as:

\begin{equation}
\theta = [ \ln(\sigma^2_L), \ln(\ell_{t,L}), \ln(\sigma^2_H), \ln(\ell_{t,H}), \rho_0, \ln(\epsilon_L), \ln(\epsilon_H), \ln(\ell_{s1,L}), \ln(\ell_{s2,L}), \dots ]^T
\end{equation}

where $\ell_t$ and $\ell_s$ represent temporal and spatial length-scales respectively, and $\epsilon$ denotes the nugget effect. A jitter term of $10^{-8}$ was added to the diagonal of the covariance matrices to ensure they remained positive definite during Cholesky decomposition.

\item {3. \textbf{Initialization and Multi-start Strategy}}
To validate the robustness of the estimation and avoid local maxima, we employed a multi-start strategy with $n=3$ independent runs:
\begin{itemize}
    \item \textbf{Warm-start:} The first run was initialized with physically informed defaults (e.g., $\ln(1.0)$ for length-scales).
    \item \textbf{Jittered Restarts:} Subsequent runs were initialized by adding Gaussian noise $\delta \sim \mathcal{N}(0, 0.10)$ to the base vector.
    \item \textbf{Sequential Initialization:} Parameters for non-stationary models were initialized using the optimal values obtained from simpler constant-parameter models.
\end{itemize}

\item{4. \textbf{Convergence Criteria}}
The optimisaton was governed by the specific tolerances and limits detailed in Table \ref{tab:optim_params}.

\begin{table}[h]
\centering
\caption{Numerical optimisaton parameters}
\label{tab:optim_params}
\begin{tabular}{@{}ll@{}}
\toprule
\textbf{Parameter} & \textbf{Value} \\ \midrule
Algorithm & Quasi-Newton (BFGS) \\
Optimality Tolerance ($\epsilon$) & $1 \times 10^{-3}$ \\
Step Tolerance & $1 \times 10^{-8}$ \\
Maximum Iterations & 200 \\
Maximum Function Evaluations & 2,000 \\
Number of Restarts & 3 \\ \bottomrule
\end{tabular}
\end{table}
\end{enumerate}
\end{reply}

\begin{point}
    Typos. The header of Section 6 is misspelled as “Apprendix.” In
Section 3.1, “Synthetic” is misspelled as “Sythentic.”
\end{point}
\begin{reply}
    Done.
\end{reply}

\begin{point}
    The cross-covariance block notation in Equation 5 is dense. Explicitly
stating the dimensions of the blocks (and possibly adding a brief “shapes”
line below the equation) would improve readability.
\end{point}
\begin{reply}
    Done.
\end{reply}

\begin{point}
Please state clearly whether validation splits are random or blocked in
space/time, and whether HF stations are held out entirely or partially.
This affects interpretation of reported accuracy.
\end{point}
\begin{reply}
We agree that the validation strategy is a critical aspect of our study. We have clarified in the revised manuscript that we held out the high-fidelity (HF) stations in their entirety. 

Specifically, in Section 3.2 (\textit{Synthetic data Experiment}), we have updated the text to state:
\begin{quote}
\textit{A training fraction of 0.3 is applied to the spatial locations; thus, 12 stations (including their complete time series) are used for model fitting, while the remaining 28 stations are reserved entirely for out-of-sample testing.}
\end{quote}

Similarly, in Section 3.3.1 (\textit{Real data experiment}), we have added the following clarification:
\begin{quote}
\textit{Each trained model is then evaluated on the excluded station to assess predictive accuracy and generalization performance.}
\end{quote}

This ensures that the model is tested on its ability to generalize to completely unobserved spatial locations across the entire temporal horizon.

\end{reply}



\begin{point}
    A short algorithm box would be helpful in detailing the implementation
pipeline. This would make the implementation pipeline clearer and
significantly improve the reproducibility of the method.
\end{point}

\begin{reply}
     Done. See Algorithm 1  in section 2.2.
\end{reply}
   

\begin{point}
    For the wind-speed application, please ensure units are stated and provide
sufficient detail on HF station matching to LF grid cells (nearest
neighbour vs. interpolation), time alignment, and any specific quality control
steps used. This is crucial for both reproducibility and interpretation
of the physical results.
\end{point}

\begin{reply}
    Done. Please see Section 3.41. We always used  nearest neighbour matching.
\end{reply}

\section{Reviewer 2}

\begin{point}
    The paper should be carefully rewritten in terms of the structure. There are so many jumps from the main text and the appendix.
\end{point}


\begin{reply}
We thank the referee for this comment. In the revised version, we have substantially restructured the manuscript to improve coherence, logical flow, and separation between core methodology and supporting material.

Specifically, the revision was guided by the following principles:

\begin{enumerate}
    \item \textbf{Self-contained main text.}  
    All essential methodological components required to understand the proposed framework---including the multi-fidelity model formulation, the Vecchia-based decomposition, the likelihood reconstruction, the GLS mean-removal strategy, and the spatially varying cross-fidelity scaling---are now presented entirely within the main body of the paper (Sections~2--4). The main text no longer relies on the appendices for key definitions or algorithmic steps.

    \item \textbf{Appendices limited to supporting material.}  
    The appendices have been rewritten and restricted to supplementary content, including formal theoretical arguments (e.g., positive semidefiniteness), detailed algorithmic derivations, simulation design, and sensitivity analyses (e.g., ordering effects and sparsity patterns). These materials now strictly support, rather than interrupt, the narrative developed in the main text.

    \item \textbf{Explicit forward referencing.}  
    Where additional technical detail is useful but not essential for continuity, we now use explicit and localized forward references (e.g., ``see Appendix~A for details''), avoiding repeated back-and-forth transitions between the main text and the appendices.

    \item \textbf{Linear methodological pipeline.}  
    The methodological section has been reorganised to follow a single, coherent progression: multi-fidelity model formulation, Vecchia approximation of latent components, covariance reconstruction via the Woodbury identity, GLS mean removal, and extension to non-stationary cross-fidelity scaling. This removes earlier structural jumps where theoretical, algorithmic, and implementation aspects were interleaved.

    \item \textbf{Alignment of experiments and methodology.}  
    The experimental section has been reorganized to mirror the methodological contributions, with a dedicated validation study for the decomposed Vecchia reconstruction followed by synthetic and real-data experiments. This alignment further reduces structural fragmentation.
\end{enumerate}

Overall, the manuscript has been carefully rewritten to ensure that the main text reads as a coherent and uninterrupted exposition of the proposed framework, while the appendices now serve a clearly secondary and supportive role.

\end{reply}



\begin{point}
    In terms of format, the appendix should start from 1 instead of 7.
On, page 14, line 50, “Details regarding the specific data generation
process can be found in Appendix 6.” Appendix 6 is just a title,
with “Appendix”, and the real appendix starts from 7.1. The authors
should make the structure of appendix correct.
\end{point}

\begin{reply}
    The appendix structure has been corrected: now it goes from A to F.  
\end{reply}

\begin{point}
    The structure of the appendix is really confusing. In their appendix 7,
7.1, 7.3, 7.4 are about data generation, but 7.2 is about computation
time. There are several appendix sections about simulation, they can
consider rewrite them and put them in one section. Some descriptions
should be provided for there section 12 additional results
\end{point}

\begin{reply}
  We thank the reviewer for this helpful comment. We agree that the appendix structure in the previous version was not sufficiently clear, especially because materials related to the simulation study were split across multiple appendix sections and the computation-time figure appeared in between them. In the revised manuscript, we have reorganized the appendix to improve readability and thematic consistency.

In particular, we have grouped the simulation-related material into a single appendix section with clearly labeled subsections, so that data generation, ordering/conditioning experiments, and related simulation results now appear together in one place. We have also moved the computation-time material to a separate appendix section dedicated to computational aspects, rather than leaving it interspersed with the simulation content. In addition, we have expanded the “Additional results” section by adding introductory text and brief descriptions for each figure/table so that the purpose of that section is clear to the reader.

Overall, these revisions were made to ensure that the appendix follows a more logical progression: methodological/theoretical supplements first, simulation design and simulation-based validation second, computational diagnostics separately, and finally additional empirical results. We believe this restructuring substantially improves the clarity and navigability of the supplementary material.
\end{reply}

\begin{point}
 (a) How the Vecchia approximation is used in the spatio-temporal setup is not clear. 
What is the ordering in the Vecchia approximation? How is the neighbour structure 
determined for spatio-temporal data? \\
(b) On page 9, the description of $\mathbf{Z}_1$ and $\mathbf{Z}_{21}$ is not clear. 
It appears that each row represents a data location. \\
(c) What covariance function is used for estimation in the simulation and real-data 
experiments? Is the same spatio-temporal covariance assumed as in the data generation? \\
(d) On page 13, how is equation (11) evaluated for each point? It appears that the 
approach assumes stationarity over time. \\
(e) On page 13, the model for $y_\rho$ only enforces spatial smoothness. Given the small 
number of spatial locations (9 in the simulation and 18 in the real data), this modelling 
choice seems questionable.

\end{point}


\begin{reply}
    We thank the referee for these detailed comments. In the revised manuscript, we have 
substantially clarified the model description and implementation. Below we address 
each point in turn.

\medskip
\textbf{(a) Vecchia approximation: ordering and neighbour structure.}  
The Vecchia approximation is now explicitly described in both the methodological and 
experimental sections. In all experiments, observations are ordered using a 
\emph{space-major spatio-temporal ordering}, where data are first grouped by spatial 
location and temporal replicates are ordered sequentially within each location. 
Conditioning sets are constructed using a fixed neighbourhood size $m$, with neighbours 
selected either according to the native geometric ordering or via correlation-based 
conditioning. The effect of alternative orderings on sparsity, fill-in, and numerical 
stability is systematically investigated in Appendix~F, including a comparison of 
space-major, time-major, time-major and random space, and random orderings.

\medskip
\textbf{(b) Interpretation of $\mathbf{Z}_1$ and $\mathbf{Z}_{21}$.}  
We have clarified that $\mathbf{Z}_1$ and $\mathbf{Z}_{21}$ are binary incidence matrices 
that map latent low-fidelity process realizations to observation locations. They encode the membership  between two sets. Each row 
corresponds to an observation, and each column corresponds to a unique spatial location. 
In the revised text, we explicitly define their dimensions and explain their role in both 
nested and non-nested designs. In particular, we now introduce $N_1$ as the number of 
distinct spatial locations and clarify how it relates to $n_L$ in nested designs and to 
the union of low- and high-fidelity locations in non-nested designs.

\medskip
\textbf{(c) Covariance function in simulation and estimation.}  
Both the simulation study and the estimation procedure use the same class of covariance 
functions. Specifically, the low-fidelity latent process and the discrepancy process are 
modeled using zero-mean Gaussian processes with separable spatio-temporal covariance 
functions, defined as the product of a spatial kernel and a temporal kernel. In all main 
experiments, squared-exponential kernels are used in both space and time. This modelling 
choice is now stated explicitly in Section~2.6, and consistency between the data 
generation mechanism and the fitted models is ensured throughout the paper.

\medskip
\textbf{(d) Evaluation of equation (11) and temporal stationarity.}  
Equation~(11) defines an empirical quantity $y_\rho(\mathbf{s})$ that is computed at each 
spatial location using the temporal co-variability between low- and high-fidelity 
observations at that location. Therefore, time points with the same spatial locations have the same $\rho$. In practice, at each spatial location where both fidelities are observed,
$\rho(\bm{s})$ corresponds to the local least-squares coefficient obtained
from regressing $y_H(\bm{s},t)$ on $y_L(\bm{s},t)$ over time. This construction implicitly assumes stationarity over time. The assumption of temporal stationarity regarding the relationship between datasets is strongly supported by the empirical evidence of the work of Colombo (2024) [https://doi.org/10.1093/jrsssc/qlaf003]. Their analysis demonstrates that predictive uncertainty does not increase as the temporal gap grows—an observation that is consistent only within a stationary context. 
We have clarified the role of $\rho(s)$ in the revised manuscript. 

\medskip
\textbf{(e) Spatial GP model for $y_\rho$ with few locations.}  
We agree that, given the small number of spatial locations, the GP model for 
$\rho(\mathbf{s})$ should not be interpreted as a highly detailed nonparametric spatial 
model. The inclusion $\rho(s)$ is primarily to demonstrate the flexibility 
of the proposed framework and to explore scenarios where spatial heterogeneity in 
cross-fidelity coupling may be present, rather than as the default or recommended choice for sparse spatial designs. We further observed that the exact same spatially varying framework, when implemented without approximation, frequently results in non-positive semidefinite (non-PSD) covariance matrices due to numerical instability. Consequently, the Vecchia approximation serves a dual purpose: it not only facilitates the analysis of large datasets but also acts as a general regularizer, enhancing the numerical robustness of the model.

\medskip
Overall, we believe these revisions substantially improve the clarity of the model 
description and address the referee’s concerns.
\end{reply}

\begin{point}
    (a) How are $n_\delta = 100, 300,$ and $500$ constructed? \\
(b) The synthetic data setup is too simple. Why are only 9 spatial locations used? 
Even the real-data experiment with 18 locations is not convincing for a spatio-temporal 
model using Vecchia approximation. The authors also state that “modelling one month of 
data is not much different from modelling two months.” It would be more convincing to 
consider setups with a larger number of spatial locations. \\
(c) On page 21, Table 3, it appears that WMFgpgp performs reasonably well, yet in Figure 2 
only GP-3D is included. The figure should also include WMFgpgp.
\end{point}


\begin{reply}
    We thank the referee for these comments, which prompted a substantial redesign of the 
simulation study. In the revised manuscript, the synthetic-data experiment has been 
reconstructed to address both the scale and interpretability concerns raised.

\medskip
\textbf{(a) Construction of $n_\delta$.}  
In the revised simulation, we no longer rely on the earlier formulation based on 
$n_\delta = 100, 300,$ and $500$. Instead, the synthetic experiment is now defined on an 
explicit spatio-temporal grid with $n_{\text{space}} \times n_{\text{space}}$ spatial 
locations and $n_{\text{time}}$ temporal replicates. The total number of latent 
discrepancy observations $n_\delta$ is therefore determined directly by the grid size, 
ensuring that the dimensionality of the discrepancy process is transparent and tied to 
the underlying spatio-temporal structure rather than introduced ad hoc.

\medskip
\textbf{(b) Size and realism of the synthetic setup.}  
We agree that the original synthetic experiment was too simplistic and did not adequately reflect the setting in which Vecchia approximations are typically most relevant. In response, we completely redesigned the simulation study. The revised experiment is now based on a regular 
6×6 spatial grid (36 locations) observed over 10 time points, yielding a total of 360 spatio-temporal observations. At the same time, we deliberately kept the overall sample size moderate so that exact Gaussian process inference remains feasible. This allows us to make a direct comparison with the non-approximated model, while still working at a scale large enough to meaningfully evaluate spatio-temporal dependence and cross-fidelity structure.

The statement regarding the limited distinction between modelling one versus two months has been removed to avoid confusion. As we understand these might be peculiarities specific to this experiments.
Though, for this dataset  the statement is substantially correct, since as we already said at point 0.17 the process temporally stationary.
Regarding the size of the real-data experiment, we consider it appropriate given that it involves approximately 26,000 observations (18 stations × 744 hourly measurements). A non-approximated multi-fidelity Gaussian process model is unable to perform even a single covariance matrix inversion at this scale without exhausting available memory resources. The revised setup focuses instead on spatial 
extrapolation by holding out entire stations, which directly reflects the motivating use 
case of sparse high-fidelity monitoring networks. While the real-data experiment remains 
limited to 18 stations due to data availability, the synthetic experiment now serves as 
the primary controlled validation of the spatio-temporal modelling assumptions.

\medskip
\textbf{(c) Model comparison and figure consistency.}  
Following the restructuring of the simulation study, we have also simplified the set of 
competing models. Several ambiguous or less interpretable variants, including WMFgpgp, 
have been removed from the revised analysis. The simulation results now focus on a 
clearly defined and exhaustive set of benchmark Gaussian process models (GP-L, GP-3D, 
GP-4D) and their comparison with the proposed multi-fidelity framework. As a result, the 
figures and tables are now fully consistent, and all models discussed in the text are 
represented in the corresponding visualizations.

\medskip
Overall, the simulation section has been rewritten from the ground up to address the 
referee’s concerns. The revised experiments use a larger and more interpretable spatial 
domain, remove ambiguous model variants, and provide a clearer and more convincing 
assessment of the proposed methodology.
\end{reply}


\begin{point}
citation format on page 3, line 16. In the parentheses, the citation format is incorrect.
\end{point}

\begin{reply}
Not sure which.
\end{reply}

\begin{point}
    There are several lower case “table”. Usually, like
    “Figure”, it should
be “Table”, especially some are the beginning of sentences, such as
“table 1” on page 5, “table 2” on page 18, “table 3” on page 19, “table
3” on page 21
\end{point}

\begin{reply}
    Done.
\end{reply}

\begin{point}
    

On page 11 line 37, N1 seems to be a number of count, but “the
unique location” is used. It should be clear here.

\end{point}

\begin{reply}
    \noindent
We thank the referee for pointing out this imprecision :  
\textit{The only term not previously defined is $N_1$, which denotes the number of distinct spatial locations supporting the latent low-fidelity process.
In a nested design, every high-fidelity location is also observed at low fidelity, so no HF-only locations exist and $N_1 = n_L$.
In a non-nested design, some high-fidelity observations may be available at locations where no low-fidelity data are observed.
In this case, $N_1$ is equal to the total number of unique spatial locations obtained by taking the union of locations where only $y_L$ is observed, locations where only $y_H$ is observed, and locations where both $y_L$ and $y_H$ are observed.}
\end{reply}

\begin{point}
    On page 12, the equation of K, what is Σw2? should it be Σwδ?
\end{point}

\begin{reply}
Done.
\end{reply}





\end{document}